\subsection{在渐近展开上的应用}

欧拉–麦克劳林公式使我们对 V, p.~238 至 242 处理的问题——即求部分和 \(s_n = \sum_{m=0}^{n} g(m)\) （相应地，余和 \(r_n = \sum_{m=n+1}^{\infty} g(m)\) ）的渐近展开——能够（在最重要的情形）给出更完全的解，这里 \(g\) 是数量函数，\(> 0\) ，单调，定义在 \([0, +\infty[\) 上。

我们只限于 g 是 (H) 函数（V, p.~252）、相对于 \(e^{x}\) 的阶为 0 的情形；换句话说，有关系 \(g' \ll g\) ；由此推出，\(g^{(k+1)} \ll g^{(k)}\) （对每个整数 k \textgreater{} 0 ，只要各导数 \(g^{(h)}\) （阶 \(h \leqslant k\) 的）都不与常数等价；V, p.~232, prop. 7）。设 p 是这样一个整数：各导数 \(g^{(h)}\) （阶 \(h \leqslant 2p\) 的）都不与常数等价。先设以 \(g(n)\) 为通项的级数有无穷和，并区分几种情形：

\(1^{\circ}\left|g^{(2p - 1)}(n)\right|\) 趋于 \(+\infty\) （随 \(n\) ）；由假设，对 \(\left|g^{(2k - 1)}(n)\right|\) （\(1\leqslant k\leqslant p\) ）亦然；此外，由于 \(g^{(2p + 1)}\) 在 \(+\infty\) 的邻域上单调，VI, p.~289 的公式 (4) 给出 \(\mathrm{T}_p(0,n) = O\big(g^{(2p)}(n + 1)\big) = o\big(g^{(2p - 1)}(n + 1)\big)\) ；对 \(x = 0\) 应用欧拉–麦克劳林公式，就得到

\[
\begin{array}{l} s _ {n} = \sum_ {m = 0} ^ {n} g (m) = \int_ {0} ^ {n + 1} g (t) d t - \frac {1}{2} g (n + 1) \\ \qquad + \sum_ {k = 1} ^ {p} \frac {b _ {2 k}}{(2 k) !} g ^ {(2 k - 1)} (n + 1) + o \big (g ^ {(2 p - 1)} (n + 1) \big) \end{array}
\]

这个和中的每一项相对于前一项都是可忽略的；把每一项相对于一个比较尺度 E 展开，就得到 \(s_{n}\) 的一个渐近展开。

\(2^{\circ}\) 现在设对某个指标 \(q\) （满足 \(1 \leqslant q \leqslant p\) 的），\(\left|g^{(2q-1)}(n)\right|\) 趋于 \(+\infty\) （随 \(n\) ），但 \(g^{(2k-1)}(n)\) （对 \(k > q\) ）趋于 0 。由于 \(g^{(2p+1)}\) 在 \(+\infty\) 的邻域上单调，积分 \(\int_0^\infty \left|g^{(2p+1)}(u)\right| du\) 收敛，于是可以写成

\[
\begin{array}{l} s _ {n} = \sum_ {m = 0} ^ {n} g (m) = \int_ {0} ^ {n + 1} g (t) d t - \frac {1}{2} g (n + 1) + \sum_ {k = 1} ^ {q} \frac {b _ {2 k}}{(2 k) !} g ^ {(2 k - 1)} (n + 1) + C \\ \qquad + \sum_ {k = q + 1} ^ {p} \frac {b _ {2 k}}{(2 k) !} g ^ {(2 k - 1)} (n + 1) + o \big (g ^ {(2 p - 1)} (n + 1) \big) \end{array}
\]

其中 C 是常数：因为

\[
\int_ {n + 1} ^ {\infty} \left| g ^ {(2 p + 1)} (u) \right| d u = O \big (g ^ {(2 p)} (n + 1) \big) = o \big (g ^ {(2 p - 1)} (n + 1) \big).
\]

当 \(g(n)\) 本身趋于 0 时同一公式仍成立。最后，当以 \(g(n)\) 为通项的级数收敛时，对余和 \(r_n = \sum_{m=n+1}^{\infty} g(m)\) 有展开

\[
\begin{array}{l} r _ {n} = \sum_ {m = n + 1} ^ {\infty} g (m) = \int_ {n + 1} ^ {\infty} g (t) d t + \frac {1}{2} g (n + 1) \\ \qquad - \sum_ {k = 1} ^ {p} \frac {b _ {2 k}}{(2 k) !} g ^ {(2 k - 1)} (n + 1) + o \big (g ^ {(2 p - 1)} (n + 1) \big). \end{array}
\]

\setcounter{exercisegroup}{0}
\refstepcounter{exercisegroup}
